\documentclass[11pt,a4paper]{article}

\usepackage[T1]{fontenc}
\usepackage[utf8]{inputenc}
\usepackage{lmodern}
\usepackage[a4paper,margin=2.5cm]{geometry}
\usepackage{setspace}
\usepackage{microtype}
\usepackage{titlesec}
\usepackage[hidelinks]{hyperref}

\setlength{\parindent}{0pt}
\setlength{\parskip}{0.8em}
\onehalfspacing

\begin{document}

\begin{center}
	{\Large\bfseries Design of Intelligent Cycling Ecosystems: Optimization of Infrastructure, Adaptive Navigation, and Immersive Simulation for Sustainable and Safe Urban Mobility}
\end{center}

\vspace{0.6em}

As urban traffic volumes continue to grow, the safety and travel efficiency of cyclists, as vulnerable road users, depend heavily on effective cooperation with motor vehicles, connected and autonomous vehicles (CAVs), and road infrastructure. However, existing research on fine-grained cyclist intention recognition still struggles to capture the heterogeneity and uncertainty of individual behavior, while interaction decision-making between cyclists and CAVs is often modeled under simplified game-theoretic assumptions, lacking interpretable cooperative decision-making mechanisms validated with real human behavioral data.

To address these gaps, this project aims to develop a cyclist–vehicle cooperative interaction model for urban intersection scenarios that combines prediction accuracy, decision robustness, and interpretability, built upon a three-stage, closed-loop research framework.

In the first stage, deep neural networks are used to recognize and predict fine-grained cyclist intentions and real-time trajectories, which serve as inputs to the downstream decision model.

In the second stage, multi-agent reinforcement learning is employed to model the dynamic interaction and negotiation among cyclists, CAVs, and road infrastructure, while explainable artificial intelligence methods are introduced to reveal the basis of model decisions and enhance their transparency and trustworthiness.

In the third stage, an immersive simulation platform integrating SUMO with virtual reality is developed, in which real cyclists interact with simulated CAVs under controlled conditions while multimodal data—including gestures, physiological signals, and audiovisual behavior—are synchronously collected to evaluate the cognitive workload and perceived safety associated with model decisions at the level of real human perception. The resulting human-factor data are then fed back into the first two stages to refine the prediction models and decision strategies, forming a closed loop linking prediction, decision-making, validation, and optimization.

Overall, this project is expected to reduce collision risks for cyclists at urban intersections and improve traffic efficiency in mixed-traffic environments.

\end{document}